% ---- score2 standard preamble (from: python3 score2_tools.py preamble) ----
\documentclass[10pt]{extarticle}
\usepackage[a4paper,margin=1.5cm]{geometry}
\usepackage[T1]{fontenc}
\usepackage{amsmath,amssymb}
\usepackage{graphicx}
\usepackage[export]{adjustbox}
\usepackage{multicol}
\usepackage{enumitem}
\usepackage{booktabs,array,tabularx}
\setlength{\columnsep}{0.7cm}
\setlength{\parindent}{0pt}
\setlength{\parskip}{3pt}
\pagestyle{empty}
\makeatletter
% figure registry (figdefs.tex is generated by the extract step)
\newcommand{\qfigdef}[4]{\@namedef{qf@file@#1}{#2}\@namedef{qf@w@#1}{#3}\@namedef{qf@wide@#1}{#4}}
\newcommand{\qf@check}[1]{\@ifundefined{qf@file@#1}{\PackageError{score2}{Unknown figure id #1}{}}{}}
\newcommand{\qf@img}[4][]{\includegraphics[width=#2\dimexpr\@nameuse{qf@w@#4}pt\relax,max width=#3,max height=0.45\textheight,#1]{\@nameuse{qf@file@#4}}}
% Full-width questions: finalize wraps every question that contains a wide figure in \qfullbegin...\qfullend,
% so the WHOLE question (stem, figures, options) is set across the page in reading order, then the two columns
% resume. The question's head - its number, stem and first figure (\qfullsplit marks the end of the head) -
% is kept on one page, so a number is never left at a page foot with its figure overleaf; the rest may break
% across pages normally, so no big gaps are left. Without a \qfullsplit, the whole question is kept together
% if it is under 45% of a page. Never write these yourself.
\newif\ifqf@full
\newif\ifqf@boxed
\newsavebox{\qf@box}
\newcommand{\qf@open}{\global\qf@fulltrue\par\global\qf@boxedtrue
  \setbox\qf@box\vbox\bgroup\hsize\textwidth\linewidth\textwidth}
\newcommand{\qf@emit}[1]{\par\egroup\global\qf@boxedfalse
  \ifdim\dimexpr\ht\qf@box+\dp\qf@box\relax<#1\textheight\noindent\usebox{\qf@box}\par\else\unvbox\qf@box\fi}
\newcommand{\qfullsplit}{\ifqf@boxed\qf@emit{0.85}\fi}
\newcommand{\qf@close}{\ifqf@boxed\qf@emit{0.45}\fi\global\qf@fullfalse}
\newcommand{\qfullbegin}{\end{multicols}\qf@open}
\newcommand{\qfullend}{\qf@close\begin{multicols}{2}}
\newcommand{\qfullstart}{\qf@open}
\newcommand{\qfulldone}{\qf@close}
% \qfig[scale]{id}: block figure, centred, directly where it is written (size set by the tool).
% Fallback if a wide figure is not inside a full-width question: break out of the columns in place - no floats.
\newcommand{\qfig}[2][1]{\qf@check{#2}%
  \ifqf@full\qf@blk{#1}{#2}\else
  \ifnum\@nameuse{qf@wide@#2}=1 \expandafter\@firstoftwo\else\expandafter\@secondoftwo\fi
  {\end{multicols}\begin{center}\qf@img{#1}{\textwidth}{#2}\end{center}\begin{multicols}{2}}%
  {\qf@blk{#1}{#2}}\fi}
\newcommand{\qf@blk}[2]{\par\smallskip{\centering\qf@img{#1}{\linewidth}{#2}\par}\smallskip}
% \qfigi[scale]{id}: inline figure (inside option lists / tables / minipages); never breaks columns.
\newcommand{\qfigi}[2][1]{\qf@check{#2}\qf@img[valign=c]{#1}{\linewidth}{#2}}
\makeatother
% structure macros - lint reads these to check completeness and figure placement
\newcommand{\papersec}[1]{\par\bigskip{\centering\Large\bfseries PAPER-#1\par}\medskip}
\newcommand{\partsec}[2]{\par\medskip\noindent{\large\bfseries PART-#1: #2}\par\smallskip}
\newcommand{\sect}[1]{\par\smallskip\noindent\textbf{#1}\par}
\newcommand{\qn}[1]{\par\smallskip\noindent\textbf{#1.}~}
\newcommand{\paraq}[2]{\par\smallskip\noindent\textbf{Paragraph for Questions #1 and #2}\par}
\newenvironment{qoptions}{\begin{itemize}[leftmargin=1.8em,itemsep=1pt,topsep=2pt,label={}]}{\end{itemize}}
\newcommand{\opt}[2]{\item[(#1)] #2}
% figure registry (generated by score2_tools.py finalize)
\qfigdef{p1_phy_s1_q01_1}{Score-II - Test 20 (23 Feb 2024)-img/p1_phy_s1_q01_1.png}{109.4}{0}
\qfigdef{p1_phy_s1_q02_1}{Score-II - Test 20 (23 Feb 2024)-img/p1_phy_s1_q02_1.png}{113.6}{0}
\qfigdef{p1_phy_s1_q07_1}{Score-II - Test 20 (23 Feb 2024)-img/p1_phy_s1_q07_1.png}{56.6}{0}
\qfigdef{p1_phy_s1_q07_2}{Score-II - Test 20 (23 Feb 2024)-img/p1_phy_s1_q07_2.png}{59.2}{0}
\qfigdef{p1_phy_s1_q07_3}{Score-II - Test 20 (23 Feb 2024)-img/p1_phy_s1_q07_3.png}{65.6}{0}
\qfigdef{p1_phy_s1_q07_4}{Score-II - Test 20 (23 Feb 2024)-img/p1_phy_s1_q07_4.png}{56.5}{0}
\qfigdef{p1_phy_s1_q08_1}{Score-II - Test 20 (23 Feb 2024)-img/p1_phy_s1_q08_1.png}{108.0}{0}
\qfigdef{p1_phy_s1_q09_1}{Score-II - Test 20 (23 Feb 2024)-img/p1_phy_s1_q09_1.png}{113.2}{0}
\qfigdef{p1_phy_s1_q10_1}{Score-II - Test 20 (23 Feb 2024)-img/p1_phy_s1_q10_1.png}{89.4}{0}
\qfigdef{p1_phy_s1_q11_1}{Score-II - Test 20 (23 Feb 2024)-img/p1_phy_s1_q11_1.png}{97.2}{0}
\qfigdef{p1_phy_s1_q12_1}{Score-II - Test 20 (23 Feb 2024)-img/p1_phy_s1_q12_1.png}{131.3}{0}
\qfigdef{p1_phy_s1_q16_1}{Score-II - Test 20 (23 Feb 2024)-img/p1_phy_s1_q16_1.png}{97.3}{0}
\qfigdef{p1_phy_s1_q17_1}{Score-II - Test 20 (23 Feb 2024)-img/p1_phy_s1_q17_1.png}{128.8}{0}
\qfigdef{p1_phy_s1_q19_1}{Score-II - Test 20 (23 Feb 2024)-img/p1_phy_s1_q19_1.png}{57.0}{0}
\qfigdef{p1_phy_s2_q01_1}{Score-II - Test 20 (23 Feb 2024)-img/p1_phy_s2_q01_1.png}{87.9}{0}
\qfigdef{p1_phy_s2_q03_1}{Score-II - Test 20 (23 Feb 2024)-img/p1_phy_s2_q03_1.png}{157.7}{0}
\qfigdef{p1_phy_s2_q04_1}{Score-II - Test 20 (23 Feb 2024)-img/p1_phy_s2_q04_1.png}{110.4}{0}
\qfigdef{p1_phy_s2_q05_1}{Score-II - Test 20 (23 Feb 2024)-img/p1_phy_s2_q05_1.png}{86.9}{0}
\qfigdef{p1_phy_s2_q06_1}{Score-II - Test 20 (23 Feb 2024)-img/p1_phy_s2_q06_1.png}{115.9}{0}
\qfigdef{p1_phy_s2_q07_1}{Score-II - Test 20 (23 Feb 2024)-img/p1_phy_s2_q07_1.png}{235.9}{0}
\qfigdef{p1_phy_s2_q09_1}{Score-II - Test 20 (23 Feb 2024)-img/p1_phy_s2_q09_1.png}{59.0}{0}
\qfigdef{p1_phy_s2_q10_1}{Score-II - Test 20 (23 Feb 2024)-img/p1_phy_s2_q10_1.png}{166.5}{0}
\qfigdef{p1_chem_s1_q08_1}{Score-II - Test 20 (23 Feb 2024)-img/p1_chem_s1_q08_1.png}{21.8}{0}
\qfigdef{p1_chem_s1_q08_2}{Score-II - Test 20 (23 Feb 2024)-img/p1_chem_s1_q08_2.png}{22.0}{0}
\qfigdef{p1_chem_s1_q08_3}{Score-II - Test 20 (23 Feb 2024)-img/p1_chem_s1_q08_3.png}{49.2}{0}
\qfigdef{p1_chem_s1_q09_1}{Score-II - Test 20 (23 Feb 2024)-img/p1_chem_s1_q09_1.png}{74.2}{0}
\qfigdef{p1_chem_s1_q09_2}{Score-II - Test 20 (23 Feb 2024)-img/p1_chem_s1_q09_2.png}{69.3}{0}
\qfigdef{p1_chem_s1_q09_3}{Score-II - Test 20 (23 Feb 2024)-img/p1_chem_s1_q09_3.png}{95.5}{0}
\qfigdef{p1_chem_s1_q09_4}{Score-II - Test 20 (23 Feb 2024)-img/p1_chem_s1_q09_4.png}{93.8}{0}
\qfigdef{p1_chem_s1_q09_5}{Score-II - Test 20 (23 Feb 2024)-img/p1_chem_s1_q09_5.png}{72.0}{0}
\qfigdef{p1_chem_s1_q12_1}{Score-II - Test 20 (23 Feb 2024)-img/p1_chem_s1_q12_1.png}{63.4}{0}
\qfigdef{p1_chem_s2_q04_1}{Score-II - Test 20 (23 Feb 2024)-img/p1_chem_s2_q04_1.png}{146.2}{0}
\qfigdef{p1_chem_s2_q06_1}{Score-II - Test 20 (23 Feb 2024)-img/p1_chem_s2_q06_1.png}{166.5}{0}
\qfigdef{p1_chem_s2_q06_2}{Score-II - Test 20 (23 Feb 2024)-img/p1_chem_s2_q06_2.png}{166.0}{0}
\qfigdef{p1_chem_s2_q06_3}{Score-II - Test 20 (23 Feb 2024)-img/p1_chem_s2_q06_3.png}{109.9}{0}
\qfigdef{p1_chem_s2_q06_4}{Score-II - Test 20 (23 Feb 2024)-img/p1_chem_s2_q06_4.png}{103.1}{0}
\qfigdef{p2_phy_s1_q03_1}{Score-II - Test 20 (23 Feb 2024)-img/p2_phy_s1_q03_1.png}{74.7}{0}
\qfigdef{p2_phy_s2_q05_1}{Score-II - Test 20 (23 Feb 2024)-img/p2_phy_s2_q05_1.png}{79.3}{0}
\qfigdef{p2_phy_s2_q06_1}{Score-II - Test 20 (23 Feb 2024)-img/p2_phy_s2_q06_1.png}{167.3}{0}
\qfigdef{p2_phy_s2_q07_1}{Score-II - Test 20 (23 Feb 2024)-img/p2_phy_s2_q07_1.png}{72.5}{0}
\qfigdef{p2_phy_s3_para3_4_1}{Score-II - Test 20 (23 Feb 2024)-img/p2_phy_s3_para3_4_1.png}{108.5}{0}
\qfigdef{p2_phy_s4_q05_1}{Score-II - Test 20 (23 Feb 2024)-img/p2_phy_s4_q05_1.png}{112.3}{0}
\qfigdef{p2_phy_s4_q07_1}{Score-II - Test 20 (23 Feb 2024)-img/p2_phy_s4_q07_1.png}{49.5}{0}
\qfigdef{p2_phy_s4_q08_1}{Score-II - Test 20 (23 Feb 2024)-img/p2_phy_s4_q08_1.png}{150.4}{0}
\qfigdef{p2_phy_s4_q09_1}{Score-II - Test 20 (23 Feb 2024)-img/p2_phy_s4_q09_1.png}{103.1}{0}
\qfigdef{p2_chem_s1_q03_1}{Score-II - Test 20 (23 Feb 2024)-img/p2_chem_s1_q03_1.png}{47.2}{0}
\qfigdef{p2_chem_s1_q03_2}{Score-II - Test 20 (23 Feb 2024)-img/p2_chem_s1_q03_2.png}{45.8}{0}
\qfigdef{p2_chem_s1_q03_3}{Score-II - Test 20 (23 Feb 2024)-img/p2_chem_s1_q03_3.png}{45.6}{0}
\qfigdef{p2_chem_s1_q03_4}{Score-II - Test 20 (23 Feb 2024)-img/p2_chem_s1_q03_4.png}{45.5}{0}
\qfigdef{p2_chem_s2_q05_1}{Score-II - Test 20 (23 Feb 2024)-img/p2_chem_s2_q05_1.png}{45.1}{0}
\qfigdef{p2_chem_s2_q06_1}{Score-II - Test 20 (23 Feb 2024)-img/p2_chem_s2_q06_1.png}{179.9}{0}
\qfigdef{p2_chem_s2_q06_2}{Score-II - Test 20 (23 Feb 2024)-img/p2_chem_s2_q06_2.png}{179.9}{0}
\qfigdef{p2_chem_s2_q06_3}{Score-II - Test 20 (23 Feb 2024)-img/p2_chem_s2_q06_3.png}{269.1}{1}
\qfigdef{p2_chem_s2_q06_4}{Score-II - Test 20 (23 Feb 2024)-img/p2_chem_s2_q06_4.png}{190.0}{0}
\qfigdef{p2_chem_s4_q07_1}{Score-II - Test 20 (23 Feb 2024)-img/p2_chem_s4_q07_1.png}{173.1}{0}
\qfigdef{p2_chem_s4_q07_2}{Score-II - Test 20 (23 Feb 2024)-img/p2_chem_s4_q07_2.png}{322.1}{1}
\qfigdef{p2_chem_s4_q08_1}{Score-II - Test 20 (23 Feb 2024)-img/p2_chem_s4_q08_1.png}{316.8}{1}
\qfigdef{p2_chem_s4_q08_2}{Score-II - Test 20 (23 Feb 2024)-img/p2_chem_s4_q08_2.png}{232.9}{0}
\qfigdef{p2_chem_s4_q10_1}{Score-II - Test 20 (23 Feb 2024)-img/p2_chem_s4_q10_1.png}{120.0}{0}
\qfigdef{p2_math_s4_q10_1}{Score-II - Test 20 (23 Feb 2024)-img/p2_math_s4_q10_1.png}{67.6}{0}
% ---- body skeleton ----
\begin{document}
\begin{center}{\Large\bfseries JEE (Main+Advanced) Score-II --- Test 20 (23 Feb 2024)}\end{center}
\begin{multicols}{2}
\raggedcolumns
\papersec{1}
\partsec{1}{PHYSICS}
\sect{Section-I}

\qn{1} A capacitor C is connected to a conducting rod of length $\ell$. The loop starts entering a region having transverse magnetic field $B$ as shown in diagram. When the loop is entering the field, what will be the charge developed on the capacitor ?
\qfig{p1_phy_s1_q01_1}
\begin{qoptions}
\opt{A}{0}
\opt{B}{$B\ell vC$}
\opt{C}{$\dfrac{B\ell vC}{2}$}
\opt{D}{$\dfrac{B\ell vC}{3}$}
\end{qoptions}

\qn{2} In the figure shown let $i_1$ be current from battery just after the switch is closed and $i_2$ be current at steady state. Find $i_1/i_2$.
\qfig{p1_phy_s1_q02_1}
\begin{qoptions}
\opt{A}{1}
\opt{B}{0.8}
\opt{C}{1.25}
\opt{D}{None}
\end{qoptions}

\qn{3} \textbf{Assertion (A) :} The flux of magnetic field through any complete closed surface is always zero but flux of electrostatic field through any closed surface may or may not be zero.\\
\textbf{Reason (R) :} The line integral of magnetic field through any closed loop may or may not be zero but line integral of electrostatic field through any closed loop is always zero.
\begin{qoptions}
\opt{A}{(A) is True, (R) is True; (R) is a correct explanation for (A).}
\opt{B}{(A) is True, (R) is True; (R) is NOT a correct explanation for (A).}
\opt{C}{(A) is True, (R) is False.}
\opt{D}{(A) is False, (R) is True.}
\end{qoptions}

\qn{4} A charged capacitor discharges through a resistance $R$ with time constant $\tau$. The two are now placed in series across an AC source of angular frequency $\omega=\dfrac1\tau$. The impedance of the circuit will be-
\begin{qoptions}
\opt{A}{$\dfrac{R}{\sqrt2}$}
\opt{B}{$R$}
\opt{C}{$\sqrt2R$}
\opt{D}{$2R$}
\end{qoptions}

\qn{5} A proton and deuteron are both accelerated through the same potential difference and enter in a magnetic field perpendicular to the direction of the field. If the deuteron follows a path of radius $R$, assuming the neutron and proton masses are nearly equal, the radius of the proton's path will be :-
\begin{qoptions}
\opt{A}{$\sqrt2R$}
\opt{B}{$\dfrac{R}{\sqrt2}$}
\opt{C}{$\dfrac R2$}
\opt{D}{$R$}
\end{qoptions}

\qn{6} An electron is moving along $+x$ direction with a velocity of $6\times10^6\ \mathrm{ms^{-1}}$. It enters a region of uniform electric field of 300 V/cm pointing along $+y$ direction. The magnitude and direction of the magnetic field set up in this region such that the electron keeps moving along the $x$ direction will be :-
\begin{qoptions}
\opt{A}{$5\times10^{-3}$ T, along $+z$ direction}
\opt{B}{$3\times10^{-4}$ T, along $-z$ direction}
\opt{C}{$3\times10^{-4}$ T, along $+z$ direction}
\opt{D}{$5\times10^{-3}$ T, along $-z$ direction}
\end{qoptions}

\qn{7} A negatively charged particle is revolving in a circle of radius $r$. Out of the following which one figure represents the correct directions $\vec L$ of and $\vec M$ ?\\
($\vec L$ is angular momentum of particle; $\vec M$ is magnetic moment of the particle).
\begin{qoptions}
\opt{A}{\qfigi{p1_phy_s1_q07_1}}
\opt{B}{\qfigi{p1_phy_s1_q07_2}}
\opt{C}{\qfigi{p1_phy_s1_q07_3}}
\opt{D}{\qfigi{p1_phy_s1_q07_4}}
\end{qoptions}

\qn{8} Two crossed wires each carrying current $I$ are shown in the figure. The direction of the force exerted on wire 2 by wire 1 at the point C \& D are shown by arrow :-
\qfig{p1_phy_s1_q08_1}
\begin{qoptions}
\opt{A}{Current in wire 1 flows from left to right and current in wire 2 flows from bottom to top.}
\opt{B}{Current in wire 1 flows from left to right and current in wire 2 flows from top to bottom}
\opt{C}{Current in wire 1 flows from right to left \& current in wire 2 flows from bottom to top.}
\opt{D}{It is impossible to have given force directions at point C \& D}
\end{qoptions}

\qn{9} The magnetic field in the cylindrical region shown in figure increases at a constant rate of 10.0 mT s$^{-1}$. Each side of the square loop abcd and defa has a length of 2.00 cm and a resistance of 2.00$\Omega$. Correctly match the current in the wire ad in four different situations as listed in list I with the values given in list II.
\qfig{p1_phy_s1_q09_1}
\begin{center}
\begin{tabular}{|c|p{3.4cm}|c|p{2.2cm}|}
\hline
 & \multicolumn{1}{c|}{\textbf{List-I}} & & \multicolumn{1}{c|}{\textbf{List-II}}\\
\hline
(P) & The switch S$_1$ is closed but S$_2$ is open & (1) & $5\times10^{-7}$ A, d to a\\
\hline
(Q) & S$_1$ is open but S$_2$ is closed & (2) & $5\times10^{-7}$ A, a to d\\
\hline
(R) & both S$_1$ and S$_2$ are open & (3) & 0 A\\
\hline
(S) & both S$_1$ and S$_2$ are closed & (4) & $2\times10^{-6}$ A, a to d\\
\hline
\end{tabular}
\end{center}
\begin{qoptions}
\opt{A}{(P)$\rightarrow$(2), (Q)$\rightarrow$(1), (R)$\rightarrow$(3), (S)$\rightarrow$(4)}
\opt{B}{(P)$\rightarrow$(2), (Q)$\rightarrow$(1), (R)$\rightarrow$(3), (S)$\rightarrow$(3)}
\opt{C}{(P)$\rightarrow$(1), (Q)$\rightarrow$(2), (R)$\rightarrow$(3), (S)$\rightarrow$(4)}
\opt{D}{(P)$\rightarrow$(2), (Q)$\rightarrow$(1), (R)$\rightarrow$(4), (S)$\rightarrow$(3)}
\end{qoptions}

\qn{10} A charge of $+10\mu$C and mass of 1gm is moving with velocity of $10^6$ m/s at an angle of $45^\circ$ to x \& y axis as shown. It experiences a magnetic force along $-$ve z-direction. When it is projected with $10^6$ m/s along $+$ve z-direction, it experiences a magnetic force of $10^{-2}$ N in $+$ve x-direction. The magnetic field $\vec B$ is :
\qfig{p1_phy_s1_q10_1}
\begin{qoptions}
\opt{A}{$-10^{-3}\,\mathrm T\left(\hat i+\hat j\right)$}
\opt{B}{$10^{-3}\left(\hat j+\hat k\right)$}
\opt{C}{$10^{-3}\,\hat i$}
\opt{D}{$-10^{-3}\,\mathrm T\,\hat j$}
\end{qoptions}

\qn{11} Figure shows a square current-carrying loop ABCD of side 2 m and current $I=\dfrac12$A. The magnetic moment $\vec M$ of the loop is :-
\qfig{p1_phy_s1_q11_1}
\begin{qoptions}
\opt{A}{$\left(\hat i-\sqrt3\hat k\right)$ A$-$m$^2$}
\opt{B}{$\left(\hat j-\hat k\right)$ A$-$m$^2$}
\opt{C}{$\left(\sqrt3\hat i-\hat k\right)$ A$-$m$^2$}
\opt{D}{$\left(\hat i-\hat k\right)$ A$-$m$^2$}
\end{qoptions}

\qn{12} A 4 kg block moving with 10 m/s strikes a spring of constant $\pi^2$ N/m attached to 4 kg block at rest kept on a smooth floor. The time for which rear moving block remain in contact with spring will be :
\qfig{p1_phy_s1_q12_1}
\begin{qoptions}
\opt{A}{$\sqrt2$ sec}
\opt{B}{$\dfrac1{\sqrt2}$ sec}
\opt{C}{1 sec}
\opt{D}{$\dfrac12$ sec}
\end{qoptions}

\qn{13} The position vector of a particle moving in x-y plane is given by $\vec r=(A\sin\omega t)\hat i+(A\cos\omega t)\hat j$ then motion of the particle is :
\begin{qoptions}
\opt{A}{SHM}
\opt{B}{on a circle}
\opt{C}{on a straight line}
\opt{D}{with constant acceleration}
\end{qoptions}

\qn{14} $Y=A\sin(\omega t+\phi_0)$ is the time-displacement equation of a simple harmonic motion. At $t=0$ the displacement of the particle is $Y=\dfrac A2$ and it is moving along negative x-direction. Then the initial phase angle $\phi_0$ will be :-
\begin{qoptions}
\opt{A}{$\dfrac\pi6$}
\opt{B}{$\dfrac\pi3$}
\opt{C}{$\dfrac{5\pi}6$}
\opt{D}{$\dfrac{2\pi}3$}
\end{qoptions}

\qn{15} \textbf{Statement-1 :} Magnetic field does not do any work on moving charges.\\
\textbf{Statement-2 :} Magnetic field lines always form closed loops.
\begin{qoptions}
\opt{A}{Statement-1 is true, statement-2 is true and statement-2 is correct explanation for statement-1.}
\opt{B}{Statement-1 is true, statement-2 is true and statement-2 is NOT the correct explanation for statement-1.}
\opt{C}{Statement-1 is true, statement-2 is false.}
\opt{D}{Statement-1 is false, statement-2 is true.}
\end{qoptions}

\qn{16} The x-t graph of a particle undergoing simple harmonic motion is shown below. The acceleration of the particle at $t=4/3$ s is :-
\qfig{p1_phy_s1_q16_1}
\begin{qoptions}
\opt{A}{$\dfrac{\sqrt3}{32}\pi^2$cm/s$^2$}
\opt{B}{$\dfrac{-\pi^2}{32}$cm/s$^2$}
\opt{C}{$\dfrac{\pi^2}{32}$cm/s$^2$}
\opt{D}{$-\dfrac{\sqrt3}{32}\pi^2$cm/s$^2$}
\end{qoptions}

\qn{17} A conducting loop with shape (a) is placed in uniform magnetic field $B$. Its one curved portion is turned so that the shape becomes (b). Find the correct statements(s). Given $\Omega$ be the resistance.
\qfig{p1_phy_s1_q17_1}
\begin{qoptions}
\opt{A}{Change in flux is given by $2B\pi R^2$.}
\opt{B}{Change in flux is given by $\dfrac{B\pi R^2}{2}$.}
\opt{C}{Total charge crossing any point of loop is $\dfrac{B\pi R^2}{2\Omega}$.}
\opt{D}{Total charge crossing any point of loop is $\dfrac{B\pi R^2}{\Omega}$.}
\end{qoptions}

\qn{18} A flexible wire loop in the shape of a circle has a radius that grows linearly with time. There is a magnetic field perpendicular to the plane of the loop that has a magnitude inversely proportional to the distance from the centre of loop i.e. $B(r)\propto\dfrac1r$. How does the emf $E$ vary with time $t$ ?
\begin{qoptions}
\opt{A}{$E\propto t^2$}
\opt{B}{$E\propto t$}
\opt{C}{$E\propto\sqrt t$}
\opt{D}{$E$ is constant}
\end{qoptions}

\qn{19} The ring B is coaxial with a solenoid A as shown in figure. As the switch S is closed at $t=0$, the ring B :-
\qfig{p1_phy_s1_q19_1}
\begin{qoptions}
\opt{A}{Is attracted towards A.}
\opt{B}{Is repelled by A.}
\opt{C}{Is initially repelled and then attracted.}
\opt{D}{Is initially attracted and then repelled.}
\end{qoptions}

\qn{20} A ferromagnetic material is heated above its curie temperature. Which one is a correct statement?
\begin{qoptions}
\opt{A}{Ferromagnetic domains are perfectly arranged.}
\opt{B}{Ferromagnetic material changes itself into paramagnetic material.}
\opt{C}{Ferromagnetic domains are not influenced.}
\opt{D}{Ferromagnetic material changes itself into diamagnetic material.}
\end{qoptions}

\sect{Section-II}

\qn{1} What is the amount of power delivered by the ac source in the circuit shown (in watts)?
\qfig{p1_phy_s2_q01_1}

\qn{2} A particle executes SHM along a straight line with mean position at $x=0$, period 20 sec and amplitude 5 cm. Find the shortest time taken by the particle to go from $x=4$ cm to $x=-3$ cm in seconds.

\qn{3} Find the magnetic field (in $\mu$T) at point A. (Given : $i=\dfrac{20}{\pi}$ Amp, $R=1$m)
\qfig{p1_phy_s2_q03_1}

\qn{4} In a region, magnetic field along X-axis changes with time according to the given graph. A charged particle of mass 1 kg \& charge 1C is located at origin starts moving with initial velocity $\left(\dfrac2\pi\hat i+2\sqrt3\hat j\right)$ m/s. It completes one revolution in $T_0$ time. Find the displacement (in m) of particle at time $\dfrac{5T_0}{4}$.
\qfig{p1_phy_s2_q04_1}

\qn{5} A loop is formed with two springs and a conducting rod of length $\ell=1$m and mass $m=1$ kg. A uniform magnetic field $B=0.01$T is directed perpendicular to the plane of the loop. At $t=0$ the rod is released when the spring is extended by a distance 4cm. What is the maximum value of the emf (in mV) (spring constant $k=200$ N/m) ?
\qfig{p1_phy_s2_q05_1}

\qn{6} The magnetic dipole moment of a uniformly charged lamina in shape of equilateral triangle of side $a$ is $M$, when rotated about axis AB with uniform $\omega$ (see figure). If a hexagonal lamina of side $a$ and having same charge density is rotated about axis A$'$B$'$ (with same angular speed $\omega$). The dipole moment generated is $5\alpha M$. Find $\alpha$.
\qfig{p1_phy_s2_q06_1}

\qn{7} In the given figure magnetic field at the center of ring (O) is $8\sqrt2$ T. Now it is turned through a $90^\circ$ about xx$'$ axis, so that two semicircular parts are mutually perpendicular. Find the value of magnetic field (in Tesla) at centre.
\qfig{p1_phy_s2_q07_1}

\qn{8} A bar magnet is demagnetized by inserting it inside a solenoid of length 0.2 m, 100 turns, and carrying a current of 5.2 A. The coercivity (A/m) of the bar magnet is :-

\qn{9} A cell is connected between the points A and C of a circular conductor ABCD with O as centre and angle AOC $=60^\circ$. If $B_1$ and $B_2$ are the magnitudes of the magnetic fields at O due to the currents in ABC and ADC respectively, then ratio $\dfrac{B_1}{B_2}$ is :-
\qfig{p1_phy_s2_q09_1}

\qn{10} Two blocks of masses $m_1=1$kg and $m_2=2$kg are connected by a spring of constant $k=100$ N/m and placed over a plank as shown. The coefficient of friction between $m_2$ and plank is 0.5 and there is no friction between $m_1$ and plank. Initially the whole system, is at rest. The plank is now pulled to right with constant acceleration $a=2$ m/s$^2$. Find the amplitude (in cm) of oscillation of block $m_1$ relative to plank.
\qfig{p1_phy_s2_q10_1}

\partsec{2}{CHEMISTRY}
\sect{Section-I}

\qn{1} H$_2$(g) is produced industrially by electrolysis of aq. NaCl using inert electrodes. How many ampere-hours would be required for production of 20 kg hydrogen :
\begin{qoptions}
\opt{A}{$5.36\times10^5$}
\opt{B}{$1.07\times10^6$}
\opt{C}{$2.68\times10^5$}
\opt{D}{$1.93\times10^5$}
\end{qoptions}

\qn{2} Kohlrausch's law states that at :
\begin{qoptions}
\opt{A}{For strong electrolyte, molar conductance varies linearly with concentration of electrolyte}
\opt{B}{At infinite dilution, each ion makes definite contribution to molar conductance of an electrolyte whatever be the nature of the other ion of the electrolyte.}
\opt{C}{At all concentration, each ion makes definite contribution to molar conductance of an electrolyte whatever be the nature of the other ion of the electrolyte.}
\opt{D}{Molar conductance increases with decrease in concentration}
\end{qoptions}

\qn{3} Efficiency of the following cell is 84 \%.\\
A(s) + B$^{2+}$(aq.) $\rightleftharpoons$ A$^{2+}$(aq.) + B(s) ; $\Delta$H$^\mathrm{o}=-285$ kJ.\\
Then the standard EMF of the cell will be\\
\relax[Use 1F = 96500 C]
\begin{qoptions}
\opt{A}{1.20 V}
\opt{B}{2.40 V}
\opt{C}{1.10 V}
\opt{D}{1.24 V}
\end{qoptions}

\qn{4} \textbf{Statement-1 :} The voltage of mercury cell remains constant for longer period of time.\\
\textbf{Statement-2 :} It is because net cell reaction does not involve ions.
\begin{qoptions}
\opt{A}{If both the statements are true and statement-2 is the correct explanation of statement-1}
\opt{B}{If both the statements are true but statement-2 is not the correct explanation of statement-1}
\opt{C}{If statement-1 is true and statement-2 is false}
\opt{D}{If statement-1 is false and statement-2 is true}
\end{qoptions}

\qn{5} A 500 mL sample of a 0.1 M Cr$^{3+}$ is electrolyzed to get elemental Cr, with a current of 96.5 A. If the remaining [Cr$^{3+}$] is 0.04 M, then the duration of process is : -
\begin{qoptions}
\opt{A}{30 sec}
\opt{B}{90 sec}
\opt{C}{10 sec}
\opt{D}{None of these}
\end{qoptions}

\qn{6} In zinc blende stucture of compound AB if atoms along one of the body diagonal of cubic unit cells are removed then the simplest formula of resulting compound will be (A are present at corners of cube)
\begin{qoptions}
\opt{A}{AB}
\opt{B}{A$_5$B$_4$}
\opt{C}{A$_{11}$B$_8$}
\opt{D}{A$_{11}$B$_{12}$}
\end{qoptions}

\qn{7} Select correct statement -
\begin{qoptions}
\opt{A}{In semiconductors, there is no band gap.}
\opt{B}{On increasing temperature, conductivity of semiconductor decreases and conductivity of metal decreases}
\opt{C}{Ferromagnetic substance become paramagnetic on heating}
\opt{D}{Ferromagnetic substance lose their magnetism in absence of external magnetic field.}
\end{qoptions}

\qn{8} Which of the following can decolourise bromine water solution?
\begin{qoptions}
\opt{A}{\qfigi{p1_chem_s1_q08_1}}
\opt{B}{\qfigi{p1_chem_s1_q08_2}}
\opt{C}{\qfigi{p1_chem_s1_q08_3}}
\opt{D}{All of these}
\end{qoptions}

\qn{9}\par\nobreak
\qfig{p1_chem_s1_q09_1}
$\xrightarrow[\mathrm{CS_2}]{\mathrm{Br_2}}$ ; Product of the reaction is :
\begin{qoptions}
\opt{A}{\qfigi{p1_chem_s1_q09_2}}
\opt{B}{\qfigi{p1_chem_s1_q09_3}}
\opt{C}{\qfigi{p1_chem_s1_q09_4}}
\opt{D}{\qfigi{p1_chem_s1_q09_5}}
\end{qoptions}

\qn{10} Which one of the following reactions does not lead to the formation of new carbon--carbon bond(s) ?
\begin{qoptions}
\opt{A}{Claisen Schmidt reaction}
\opt{B}{Cannizaro's reaction}
\opt{C}{Perkin reaction}
\opt{D}{Aldol reaction}
\end{qoptions}

\qn{11} CH$_3$CHO and CH$_3$COCH$_3$ can not be distinguished by :-
\begin{qoptions}
\opt{A}{Fehling solution}
\opt{B}{Lucas reagent}
\opt{C}{Schiff's reagent}
\opt{D}{Tollen's reagent}
\end{qoptions}

\qn{12} In the following structure,
\qfig{p1_chem_s1_q12_1}
Anomeric carbon is
\begin{qoptions}
\opt{A}{1}
\opt{B}{2}
\opt{C}{3}
\opt{D}{4}
\end{qoptions}

\qn{13} Which of the following is monosaccharide ?
\begin{qoptions}
\opt{A}{Galactose}
\opt{B}{Sucrose}
\opt{C}{Glycogen}
\opt{D}{Cellulose}
\end{qoptions}

\qn{14} Match the list -I with list - II.
\begin{center}
\begin{tabular}{|c|l|c|l|}
\hline
 & \multicolumn{1}{c|}{\textbf{List-I}} & & \multicolumn{1}{c|}{\textbf{List-II}}\\
\hline
 & \multicolumn{1}{c|}{\textbf{(Class of Drug)}} & & \multicolumn{1}{c|}{\textbf{(Example)}}\\
\hline
(a) & Antacid & (i) & Novestrol\\
\hline
(b) & Artificial sweetener & (ii) & Cimetidine\\
\hline
(c) & Antifertility & (iii) & Valium\\
\hline
(d) & Tranquilizers & (iv) & Alitame\\
\hline
\end{tabular}
\end{center}
\begin{qoptions}
\opt{A}{(a) $-$ (ii), (b) $-$ (iv),(c) $-$ (i), (d) $-$ (iii)}
\opt{B}{(a) $-$ (iv), (b) $-$ (i),(c) $-$ (ii), (d) $-$ (iii)}
\opt{C}{(a) $-$ (iv), (b) $-$ (iii),(c) $-$ (i), (d) $-$ (ii)}
\opt{D}{(a) $-$ (ii), (b) $-$ (iv),(c) $-$ (iii), (d) $-$ (i)}
\end{qoptions}

\qn{15} Which of the following complex compound /ion obey the sidewick EAN rule :-
\begin{qoptions}
\opt{A}{K$_3$[Fe(CN)$_6$]}
\opt{B}{[Fe(H$_2$O)$_5$NO]SO$_4$}
\opt{C}{[Fe(CO)$_5$]}
\opt{D}{[NiCl$_4$]$^{2-}$}
\end{qoptions}

\qn{16} Find the number of complex in which configuration of central metal ion is $t_{2g}^6e_g^0$ :-
\begin{qoptions}
\opt{A}{[Co(NH$_3$)$_5$Cl]Cl$_2$}
\opt{B}{K$_3$[Co(NO$_2$)$_6$]}
\opt{C}{Na$_2$[Fe(CN)$_5$NO]}
\opt{D}{All}
\end{qoptions}

\qn{17} In [Ni(DMG)$_2$], which of the following statement is incorrect :-
\begin{qoptions}
\opt{A}{It is square planar and diamagnetic complex}
\opt{B}{Oxygen atoms forms 5 members chelation with Ni$^{2+}$ ion}
\opt{C}{It is optically inactive due to presence of centre of symmetry}
\opt{D}{H-bonding can take place}
\end{qoptions}

\qn{18} Which of the following compound is diamagnetic and coloured ?
\begin{qoptions}
\opt{A}{Na$_3$[FeF$_6$]}
\opt{B}{[Cu(NH$_3$)$_4$]SO$_4$}
\opt{C}{K$_2$[MnO$_4$]}
\opt{D}{K$_2$[CrO$_4$]}
\end{qoptions}

\qn{19} Exchange of ligands in complex results in colour change. In the view of the above statement which of the following is incorrect ?
\begin{qoptions}
\opt{A}{Co(H$_2$O)$_6^{+2}$ $\rightarrow$ CoCl$_4^{-2}$ \ (Pink $\rightarrow$ Blue)}
\opt{B}{K$_4$[Fe(CN)$_6$] $\rightarrow$ Fe$_4$[Fe(CN)$_6$] \ (Yellow $\rightarrow$ Blue)}
\opt{C}{K$_3$[Cu(CN)$_4$] $\rightarrow$ [Cu(NH$_3$)$_4$]SO$_4$ \ (Colorless $\rightarrow$ Blue)}
\opt{D}{[Cr(H$_2$O)$_6$]$^{+2}$ $\rightarrow$ [Cr(NH$_3$)$_6$]$^{+3}$ \ (Blue $\rightarrow$ Violet)}
\end{qoptions}

\qn{20} \textbf{Statement-I :} 1 mol aqueous solution of CoCl$_3$.5NH$_3$ gave 2 mol AgCl on treating with AgNO$_3$.\\
\textbf{Statement-II :} Solution conductivity corresponds to 1 : 1 electrolyte.
\begin{qoptions}
\opt{A}{Both Statement I and Statement II are correct.}
\opt{B}{Statement I is correct but Statement II is incorrect.}
\opt{C}{Both Statement I and Statement II are incorrect.}
\opt{D}{Statement I is incorrect but Statement II is correct.}
\end{qoptions}

\sect{Section-II}

\qn{1} During an electrolysis of conc. H$_2$SO$_4$, perdisulphuric acid (H$_2$S$_2$O$_8$) and O$_2$ form at anode in equimolar amount. The mole of H$_2$ that will form simultaneously at other electrode will be $x$ times of O$_2$. What is the value of $x$?\\
(Given : 2H$_2$SO$_4$ $\rightarrow$ H$_2$S$_2$O$_8$ + 2H$^+$ + 2e$^-$)

\qn{2} Cu(s) + Sn$^{2+}$ (0.001M) $\rightarrow$ Cu$^{2+}$ (0.01M) + Sn(s)\\
The Gibbs free energy change for the above reaction at 298 K is $x\times10^{-1}$ kJ mol$^{-1}$;\\
The value of $x$ is\underline{\hspace{3em}}. [nearest integer]\\
\relax[Given : $E^\Theta_{\mathrm{Cu^{2+}/Cu}}=0.34$ V ; $E^\Theta_{\mathrm{Sn^{2+}/Sn}}=-0.14$ V; $F=96500\mathrm{C\ mol^{-1}}$]

\qn{3} An ionic compoud $\left(\overset{+}{\mathrm X}\overset{-}{\mathrm Y}\right)$ is doped with controlled amount of (MY$_2$). Number of cation vacancies in 1 mol of XY is found to be $6.02\times10^{21}$. If atomic mass of X, Y and M are 15,10 \& 20 a.m.u respectively : (Assume volume remains unchanged)\\
(Given : $N_A=6.02\times10^{23}$) If the \% decreases in density is $x$, then the value of 10$x$ is :

\qn{4} Identify total number of compounds give positive test with tollens reagent (silver mirror test)
\par\nobreak
\qfig{p1_chem_s2_q04_1}

\qn{5} How many number of reagents can be used to convert CH$_3$CHO into respective carboxylic acid or its salt.\\
(i) Ag(NH$_3$)$_2$OH \quad (ii) CuSO$_4$(Aq.) + Roschell salt\\
(iii) H$^+$ / KMnO$_4$ \quad (iv) Cu / $\Delta$ \quad (v) LiAlH$_4$\\
(vi) K$_2$Cr$_2$O$_7$/ H$^+$ \quad (vii) B$_2$H$_6$ + THF ; H$_2$O\\
(viii) m-CPBA \quad (ix) CH$_3$MgBr ; H$_2$O

\qn{6} The compound among following which can give 2,4-DNP test (2,4-dinitrophenylhydrazine test)
\par\nobreak
\qfig{p1_chem_s2_q06_1}
\qfig{p1_chem_s2_q06_2}
\qfig{p1_chem_s2_q06_3}
\qfig{p1_chem_s2_q06_4}

\qn{7} Total number of geometrical isomers are possible for [Co(NH$_3$)$_3$(H$_2$O)$_2$(NO$_3$)]$^{2+}$ is :-

\qn{8} Among the following total number of complexes which are low spin complex :-\\
K$_2$[PtCl$_6$] , [CrO$_4$]$^{2-}$ , [Co(H$_2$O)$_6$]$^{3+}$ , Na$_2$[Zn(CN)$_4$] , [IrF$_6$]$^{3-}$ , [Cr(en)$_3$]Cl$_3$ , [Fe(NH$_3$)$_6$]$^{2+}$

\qn{9} Among the following complexes, how many have `spin only' magnetic moment of 2.83 B.M.?\\[0pt]
[Ni(CO)$_4$], [Ni(CN)$_4$]$^{2-}$, [Ni(NH$_3$)$_6$]$^{+2}$, [NiCl$_4$]$^{2-}$, [NiF$_6$]$^{2-}$, [NiF$_6$]$^{4-}$

\qn{10} In how many species the central atom is sp$^3$ hybridized\\
PCl$_4^+$, SiF$_4$, [Ni(CN)$_4$]$^{2-}$, [Cu(NH$_3$)$_4$]$^{2+}$, H$_3$O$^+$\\
Ni(CO)$_4$, SF$_6$, [CuCl$_4$]$^{2-}$, [PtCl$_4$]$^{2-}$

\partsec{3}{MATHEMATICS}
\sect{Section-I}

\qn{1} Coefficeint of $x^{-1003}$ in the binomial expansion of $\left(x-\dfrac{1}{\sqrt x}\right)^{2009}$ is
\begin{qoptions}
\opt{A}{$-2009$}
\opt{B}{2009}
\opt{C}{2010}
\opt{D}{2008}
\end{qoptions}

\qn{2} Remainder when $7^{636}$ is divided by 19 is equal to -
\begin{qoptions}
\opt{A}{1}
\opt{B}{7}
\opt{C}{12}
\opt{D}{13}
\end{qoptions}

\qn{3} \textbf{Assertion (A) :} The number of different dissimilar terms in a complete homogeneous expression of degree 11 in three variables is 286.\\
\textbf{and}\\
\textbf{Reason (R) :} The number of ways of distribution of $m$ identical objects in $n$ person such that any person may get none, one, two or all objects are ${}^{m+n-1}C_{n-1}$.
\begin{qoptions}
\opt{A}{\textbf{A} is True, \textbf{R} is True ; \textbf{R} is a correct explanation for \textbf{A}.}
\opt{B}{\textbf{A} is True, \textbf{R} is True ; \textbf{R} is not a correct explanation for \textbf{A}.}
\opt{C}{\textbf{A} is True, \textbf{R} is False.}
\opt{D}{\textbf{A} is False, \textbf{R} is True.}
\end{qoptions}

\qn{4} If the 9$^{\text{th}}$ term in the expansion of $\left(\dfrac1{x^2}+\dfrac x2\log_2x\right)^{12}$ is equal to 495, then the value of $x$ can be -
\begin{qoptions}
\opt{A}{1}
\opt{B}{an integer greater than 1}
\opt{C}{an integer less than 1}
\opt{D}{an irrational number}
\end{qoptions}

\qn{5} If the letters of the word ``nagpur'' are written in all possible ways and then are arranged as in a dictionary, then the rank of the word ``nagpur'' is :
\begin{qoptions}
\opt{A}{241}
\opt{B}{242}
\opt{C}{243}
\opt{D}{245}
\end{qoptions}

\qn{6} There are three parallel lines lying on the same plane. If 6 distinct points are taken on each of these lines, then maximum possible number of triangles with vertices at these points is
\begin{qoptions}
\opt{A}{756}
\opt{B}{396}
\opt{C}{306}
\opt{D}{486}
\end{qoptions}

\qn{7} The number of integral solutions for $a,b,c$ and $d$ of the equation $a\cdot b\cdot c\cdot d=2021$ is :
\begin{qoptions}
\opt{A}{16}
\opt{B}{32}
\opt{C}{64}
\opt{D}{128}
\end{qoptions}

\qn{8} A box contains 20 identical balls of which 5 are white and 15 black. The balls are drawn at random from the box one at a time with replacement. The probability that a white ball is drawn for the 3$^{\text{rd}}$ time on the 6$^{\text{th}}$ draw is-
\begin{qoptions}
\opt{A}{$\dfrac12$}
\opt{B}{$\dfrac{135}{2048}$}
\opt{C}{$\dfrac{135}{1024}$}
\opt{D}{$\dfrac{27}{4096}$}
\end{qoptions}

\qn{9}
\begin{center}
\begin{tabular}{|c|p{4.2cm}|c|c|}
\hline
 & \multicolumn{1}{c|}{\textbf{List-I}} & \multicolumn{2}{c|}{\textbf{List-II}}\\
\hline
(A) & If $\displaystyle\sum_{r=0}^{10}(r+2^r)\,{}^{10}C_r=10.2^a+3^b$, then $a+b$ is equal to & (I) & 7\\
\hline
(B) & If 4$^{\text{th}}$ term in the expansion of $(x+(ax)^{-1})^m$, $m\in N$ is $\dfrac52$, then $ma$ is equal to & (II) & 12\\
\hline
(C) & Let $\lambda$ be the value of numerically greatest term in the expansion of $(1+5x)^{10}$ at $x=1$, then $\dfrac{\lambda}{{}^{10}C_9\cdot5^7}$ is equal to & (III) & 19\\
\hline
(D) & If ${}^nP_r=5040\left({}^{n-1}C_r+{}^{n-1}C_{r-1}\right)$, then $r$ is equal to & (IV) & 25\\
\hline
\end{tabular}
\end{center}
\begin{qoptions}
\opt{A}{(A)$\rightarrow$(IV); (B)$\rightarrow$(II); (C)$\rightarrow$(III); (D)$\rightarrow$(I)}
\opt{B}{(A)$\rightarrow$(I); (B)$\rightarrow$(III); (C)$\rightarrow$(II); (D)$\rightarrow$(IV)}
\opt{C}{(A)$\rightarrow$(I); (B)$\rightarrow$(II); (C)$\rightarrow$(III); (D)$\rightarrow$(IV)}
\opt{D}{(A)$\rightarrow$(III); (B)$\rightarrow$(II); (C)$\rightarrow$(IV); (D)$\rightarrow$(I)}
\end{qoptions}

\qn{10} One hundred identical coins, each with probability $p$, of showing up heads are tossed once. If $0<p<1$ and the probability of heads showing on 50 coins is equal to that of heads on 51 coins, then the value of $p$ is :
\begin{qoptions}
\opt{A}{$\dfrac12$}
\opt{B}{$\dfrac{49}{101}$}
\opt{C}{$\dfrac{50}{101}$}
\opt{D}{$\dfrac{51}{101}$}
\end{qoptions}

\qn{11} For two data sets, each of size 5, the variance are given to be 4 and 5 and the corresponding means are given to be 2 and 4, respectively. The double of the variance of the combined data set is :
\begin{qoptions}
\opt{A}{13}
\opt{B}{12}
\opt{C}{11}
\opt{D}{10}
\end{qoptions}

\qn{12} If mean of set of $n$ numbers $a_1,a_2,\ldots\ldots,a_n$ is $\bar a$, then mean of numbers $a_i+2i$, $1\le i\le n$ will be-
\begin{qoptions}
\opt{A}{$\bar a+2$}
\opt{B}{$\bar a+n$}
\opt{C}{$\bar a+2n$}
\opt{D}{$\bar a+n+1$}
\end{qoptions}

\qn{13} Let $P=\dfrac{5}{\frac1{\log_2x}+\frac1{\log_3x}+\frac1{\log_4x}+\frac1{\log_5x}}$ and $(120)^P=32$, then the value of $x$ be
\begin{qoptions}
\opt{A}{1}
\opt{B}{2}
\opt{C}{3}
\opt{D}{4}
\end{qoptions}

\qn{14} Number of solution(s) of the equation $\log_{x^2+6x+8}\log_{2x^2+2x+3}(x^2-2x)=0$ is
\begin{qoptions}
\opt{A}{1}
\opt{B}{2}
\opt{C}{3}
\opt{D}{0}
\end{qoptions}

\qn{15} Consider $(a,b)$, where $a$ and $b$ are respective out comes in throwing an unbiased die twice.\\
\textbf{Statement-1 :} If $\displaystyle\lim_{x\to0}\left(\dfrac{a^x+b^x+1}{3}\right)^{\frac3x}=6$, then the probability that `$a$' is a prime number is $\dfrac12$.\\
\textbf{and}\\
\textbf{Statement-2 :} If $A$ \& $B$ are two events then probability of event $B$ when event $A$ has already happened is $P(B/A)=\dfrac{P(A\cap B)}{P(A)}$.
\begin{qoptions}
\opt{A}{Statement-1 is True, Statement-2 is True.}
\opt{B}{Statement-1 is False, Statement-2 is False.}
\opt{C}{Statement-1 is True, Statement-2 is False.}
\opt{D}{Statement-1 is False, Statement-2 is True.}
\end{qoptions}

\qn{16} Solution set of $x+\log_{10}(1+2^x)=x\log_{10}5+\log_{10}6$ is
\begin{qoptions}
\opt{A}{$\phi$}
\opt{B}{$\{1\}$}
\opt{C}{$\left\{1,\sqrt2\right\}$}
\opt{D}{$\left\{1,\sqrt2,4\right\}$}
\end{qoptions}

\qn{17} Number of solutions of $5\cos^2\theta-3\sin^2\theta+6\sin\theta\cos\theta=7$ in the interval $[0,2\pi]$ is :-
\begin{qoptions}
\opt{A}{2}
\opt{B}{4}
\opt{C}{0}
\opt{D}{None of these}
\end{qoptions}

\qn{18} If $x$ \& $y$ are solutions of equation, $3\cos x+4\sin x=2y^2-4y+7$, then which of the following is true -
\begin{qoptions}
\opt{A}{$\tan(2xy)=\dfrac34$}
\opt{B}{$\cot(2xy)=\dfrac34$}
\opt{C}{$\tan(xy)=\dfrac34$}
\opt{D}{$\cot(xy)=\dfrac34$}
\end{qoptions}

\qn{19} Set of values of $k$ if the equation $4\sin^2x+(k+9)\sin x+2(k+1)=0$ have at least one solution, is
\begin{qoptions}
\opt{A}{$[-5,5]$}
\opt{B}{$[3,5]$}
\opt{C}{$[-5,3]$}
\opt{D}{$[-3,5]$}
\end{qoptions}

\qn{20} Mr. X walk up 16 steps, going up either 1 or 2 steps at a time. There is an explosive material on the 8$^{\text{th}}$ step so he cannot step there. The number of ways in which Mr. X can go up 16 steps is
\begin{qoptions}
\opt{A}{441}
\opt{B}{420}
\opt{C}{462}
\opt{D}{400}
\end{qoptions}

\sect{Section-II}

\qn{1} Let $n$ be the coefficient of $x^6$ in $(1-x^3)^5(1+x^2)^4(1+x^4)^8$, when expanded in powers of `x'. The value of $\left[\dfrac n{19}\right]$ is (where $[.]$ denotes greatest integer function)

\qn{2} If the number of triplets $(l,m,n)$ satisfying $1\le l\le m\le n\le20$ is ${}^pC_q$, where $l,m,n,p,q\in N$, then least value of $\dfrac{p+q}{5}$ is :

\qn{3} There are 3 papers of 100 marks each in an examination, total number of ways in which a student gets 150 marks such that he gets at least 60\% in each of two papers are $N$, then $\left[\dfrac N{400}\right]$ is equal to (where $[.]$ denotes greatest integer function)\\
(Assume that marks are all in non negative integral values)

\qn{4} Words from the letters of the word ALLENITES are formed by taking all letters at a time. The probability that both L's are not together and both E's are not together is $\dfrac ab$ (where $\dfrac ab$ is written in lowest forms), then $b-a$ equals

\qn{5} Two squares of $(1\times1)$ dimension are chosen randomly from a $(n\times n)$ chessboard. If probability that chosen squares have one side in common is $\dfrac2{45}$, then $n$ is equal to

\qn{6} Mean of 100 observations is 45. It was later found that two observations 91 and 13 were incorrectly recorded as 19 and 31. The correct mean is $\lambda$, then value of $(\lambda-40.54)$ is :

\qn{7} If the mean and variance of eight numbers 3, 7, 9, 12, 13, 20, $x$ and $y$ be 10 and 25 respectively, then $x\cdot y$ is equal to\underline{\hspace{2.5em}}

\qn{8} If $x>0$, $y>0$ such that $2^{\ell ny}=y^{\ell nx}$ and $3^{\ell nx}=y^{\ell n9}$, then minimum value of $(x^2+y^2)$ is

\qn{9} The number of values of $x$ in $[0,10\pi]$ for which $(\sin x)^{15}-(\cos2x)^{19}=2$ is

\qn{10} The number of rational terms in the expansion of $\left(5^{2/3}+10^{-1/4}\right)^{20}$ is :
\papersec{2}
\partsec{1}{PHYSICS}
\sect{Section-I(i)}

\qn{1} An electron of mass `m' is accelerated through a potential difference of $V$ and then it enters a magnetic field of induction $B$ normal to the field lines. Then the radius of the circular path is
\begin{qoptions}
\opt{A}{$\sqrt{\dfrac{2eV}{m}}$}
\opt{B}{$\sqrt{\dfrac{2Vm}{eB^2}}$}
\opt{C}{$\sqrt{\dfrac{2Vm}{eB}}$}
\opt{D}{$\sqrt{\dfrac{2Vm}{e^2B}}$}
\end{qoptions}

\qn{2} There exists a uniform magnetic and electric field of magnitude 1T and 1V/m respectively along positive $y$-axis. A charged particle of mass 1 kg and of charge 1C is having velocity 1m/sec along $+x$-axis and is at origin at $t=0$. Then the co-ordinates of particle at time $t=\pi$ seconds will be :-
\begin{qoptions}
\opt{A}{$(0,1,2)$}
\opt{B}{$(0,-\pi^2/2,-2)$}
\opt{C}{$(2,\pi^2/2,2)$}
\opt{D}{$(0,\pi^2/2,2)$}
\end{qoptions}

\qn{3} A current carrying loop is placed in a uniform magnetic field pointing in negative $z$ direction. Branch PQRS is a three quarter circle, while branch PS is straight. If force on branch PS is $F$, force on branch PQR is :-
\qfig{p2_phy_s1_q03_1}
\begin{qoptions}
\opt{A}{$\sqrt2F$}
\opt{B}{$\dfrac{F}{\sqrt2}$}
\opt{C}{$\dfrac{\pi F}{\sqrt2}$}
\opt{D}{$\sqrt2\pi F$}
\end{qoptions}

\qn{4} Time period of a particle executing simple harmonic motion is 8 sec. At $t=0$ it is at the mean position. The ratio of the distance covered by the particle in the 1st second to the 2nd second is :-
\begin{qoptions}
\opt{A}{$\dfrac{1}{\sqrt2+1}$}
\opt{B}{$\sqrt2$}
\opt{C}{$\dfrac{1}{\sqrt2}$}
\opt{D}{$\sqrt2+1$}
\end{qoptions}

\sect{Section-I(ii)}

\qn{5} Consider the LC circuit below. At the time shown the current has its maximum value. At this time which of the following is/are incorrect ?
\qfig{p2_phy_s2_q05_1}
\begin{qoptions}
\opt{A}{The charge on the capacitor has its maximum value}
\opt{B}{The magnetic field inside the inductor is zero}
\opt{C}{The electric field inside the capacitor has its maximum value}
\opt{D}{The charge on the capacitor is zero}
\end{qoptions}

\qn{6} A closed conducting loop, having resistance $R$, is being rotated about an axis perpendicular to the magnetic field. Magnetic flux through the closed conducting loop is continuously changing according to the graph shown in the adjacent figure. Then, which of the following statement(s) is/are correct ?
\qfig{p2_phy_s2_q06_1}
\begin{qoptions}
\opt{A}{The electric current through the loop is minimum (zero) at $t=1$ s, 3s and 5s.}
\opt{B}{The electric current through the loop is minimum (zero) at $t=0$ s, 2s and 6s.}
\opt{C}{Total charge flown through any cross-section of a closed conducting loop between 0 and 6 s is zero}
\opt{D}{Total work done in rotating the loop in the magnetic field is zero}
\end{qoptions}

\qn{7} In the given alternating current circuit, choose the CORRECT statement(s).
\qfig{p2_phy_s2_q07_1}
\begin{qoptions}
\opt{A}{Impedance of circuit is $2\Omega$.}
\opt{B}{Power factor of circuit is $\dfrac{1}{\sqrt2}$.}
\opt{C}{Peak value of current through resistance is $50\sqrt2$A.}
\opt{D}{Average power supplied by source is 2500W.}
\end{qoptions}

\sect{Section-II(i)}

\paraq{1}{2}
In a circuit, a metal filament lamp is connected in series with a capacitor of capacitance $C\ \mu$F across a 200 V, 50 Hz supply. The power consumed by the lamp is 500 W while the voltage drop across it is 100 V. Assume that there is no inductive load in the circuit. Take rms values of the voltages. The magnitude of the phase-angle (in degrees) between the current and the supply voltage is $\varphi$. Assume, $\pi\sqrt3\approx5$.

\qn{1} The value of $C$ is \underline{\hspace{3em}}.

\qn{2} The value of $\varphi$ is \underline{\hspace{3em}}.

\paraq{3}{4}
Momentum ($P$) displacement ($x$) graph of a particle of unit mass performing simple harmonic motion along $x$-axis is shown in the figure.
\qfig{p2_phy_s3_para3_4_1}

\qn{3} Angular frequency (in rad/sec) of the simple harmonic motion is :-

\qn{4} Total energy (in Joule) of oscillation is (potential energy at mean position $=0$) :-

\sect{Section-II(ii)}

\qn{5} A uniform magnetic field of 0.06 T is inside plane of the figure. The resistance of rod is 25 $\Omega$, mass is 36 gm and it can slide freely on smooth parallel rails which are perfectly conducting. The whole system is in horizontal plane. The rod starts from rest. If terminal velocity (in m/s) is $v$ then find $\dfrac{v}{300}$.
\qfig{p2_phy_s4_q05_1}

\qn{6} A torsional pendulum is formed by attaching a wire to the center of a meter stick with a mass of 6 kg. If the resulting period is 2 second, what is the torsional constant (in SI unit) for the wire? [Take : $\pi^2=10$]

\qn{7} Consider the given system. All string, pulleys and spring are ideal, mass $M$ is in equilibrium. Maximum amplitude of oscillation of block in vertical direction is $\dfrac{\alpha Mg}{k}$. Find value of $\alpha$.
\qfig{p2_phy_s4_q07_1}

\qn{8} A non relativistic positive charge particle of charge $q$ and mass $m$ is projected perpendicular to uniform magnetic field $B$ as shown. Neglecting gravity calculate $X$-coordinate of point on screen at which the charge particle will hit : $d=\dfrac{r\sqrt3}{2}$, where $r=\dfrac{mV}{qB}$. If the $x$-coordinate is $nr$ then find $n$.
\qfig{p2_phy_s4_q08_1}

\qn{9} A particle has a charge of 4nC. When it moves with a velocity $v_1=3\times10^4$ m/s at $45^\circ$ above $y$ axis in $yz$ plane, a uniform magnetic field exerts a force of $F_1$ along $-$ve $x$ axis on it. When it moves with a velocity of $v_2=2\times10^4$ m/s along $x$ axis, the same magnetic field exerts a force $F_2$ of $16\times10^{-5}$ N along $y$ axis. What is the magnitude of the magnetic field (in T)?
\qfig{p2_phy_s4_q09_1}

\qn{10} A forced oscillator is acted upon by a force $F=F_0\sin(\omega t)$. The amplitude of the oscillator is given by $55/\sqrt{2\omega^2-36\omega+9}$. What is the resonant angular frequency ?
\partsec{2}{CHEMISTRY}
\sect{Section-I(i)}

\qn{1} A current of 0.0965 ampere is passed for 1000 seconds through 50ml of 0.1 M NaCl, using inert electrodes. The average concentration of OH$^{-}$ in the final solution is
\begin{qoptions}
\opt{A}{0.01 M}
\opt{B}{0.02 M}
\opt{C}{0.03 M}
\opt{D}{0.04 M}
\end{qoptions}

\qn{2} An element X (At. wt. $=24$) form FCC lattice. If the edge length of unit cell is $4\times10^{-8}$ cm and the observed density is $2.4\times10^{3}$ kg/m$^3$. Then the percentage occupancy of lattice point by element X is : ($N_A=6\times10^{23}$)
\begin{qoptions}
\opt{A}{96}
\opt{B}{98}
\opt{C}{99}
\opt{D}{99.9}
\end{qoptions}

\qn{3} Which of the following compound gives iodoform test with NaOH + I$_2$ :
\begin{qoptions}
\opt{A}{\qfigi{p2_chem_s1_q03_1}}
\opt{B}{\qfigi{p2_chem_s1_q03_2}}
\opt{C}{\qfigi{p2_chem_s1_q03_3}}
\opt{D}{\qfigi{p2_chem_s1_q03_4}}
\end{qoptions}

\qn{4} The correct name of $[\mathrm{Pt(NH_3)_4Cl_2}]\,[\mathrm{PtCl_4}]$ is $-$
\begin{qoptions}
\opt{A}{Tetraamminedichloroplatinum(IV) tetrachloroplatinate(II)}
\opt{B}{Dichlorotetraammineplatinium(IV) tetrachloroplatinate(II)}
\opt{C}{Tetrachloroplatinum(II) tetraammineplatinate(IV)}
\opt{D}{Tetrachloroplatinum(II) dichlorotetraamineplatinate(IV)}
\end{qoptions}

\sect{Section-I(ii)}

\qn{5} Reaction in which aldehyde is obtained as product :
\begin{qoptions}
\opt{A}{$\mathrm{HC{\equiv}CH}\ \xrightarrow[\mathrm{dil\ H_2SO_4}]{\mathrm{HgSO_4}}$}
\opt{B}{$\mathrm{CH_3{-}C{\equiv}CH}\ \xrightarrow[\mathrm{(ii)H_2O_2/OH^-}]{\mathrm{(i)BH_3/THF}}$}
\opt{C}{\qfigi{p2_chem_s2_q05_1}\ $\xrightarrow[\mathrm{BaSO_4}]{\mathrm{H_2/Pd}}$}
\opt{D}{$\mathrm{CH_3{-}CH_2{-}OH}\ \xrightarrow[\mathrm{300^oC}]{\mathrm{Cu}}$}
\end{qoptions}

\qfullbegin
\qn{6} Which reaction sequence(s) represent correct major product :
\begin{qoptions}
\opt{A}{\qfigi{p2_chem_s2_q06_1}}
\opt{B}{\qfigi{p2_chem_s2_q06_2}}
\opt{C}{\qfigi{p2_chem_s2_q06_3}}
\opt{D}{\qfigi{p2_chem_s2_q06_4}}
\end{qoptions}
\qfullend

\qn{7} Which of the following ores is(are) concentrated by froth floatation?
\begin{qoptions}
\opt{A}{haematite}
\opt{B}{galena}
\opt{C}{copper pyrite}
\opt{D}{bauxite}
\end{qoptions}

\sect{Section-II(i)}

\paraq{1}{2}
For reversible Galvanic cell, thermodynamic parameters are given as :\\
$(\Delta G)_{T,P}=\Delta H-T\Delta S$\\
$(\Delta G)_{T,P}=-nFE_{cell}$\\
$\Delta S=nF\left[\dfrac{dE_{cell}}{dt}\right]_P$\\
$W_{useful}=-(\Delta G)_{T,P}$\\
For given cell Pt $|$ H$_2$ $|$ HCl $|$ AgCl $|$ Ag, variation of e.m.f. of the cell with temperature is given as\\
$E_{cell}$ (in volt) $=0.4-[2\times10^{-6}(T)^2]$\\
($T$ = temperature in Kelvin). (Use : 1F = 96500 C)

\qn{1} For above cell reaction, entropy change per mole of Ag (in J/K) at 27$^\circ$C will be\\
(Use : 1F = 96500 C)

\qn{2} For given cell Pt $|$ H$_2$ $|$ HCl $|$ AgCl $|$ Ag, variation of e.m.f. of the cell with temperature is given as\\
$E_{cell}$ (in volt) $=0.4-[2\times10^{-6}(T)^2]$\\
($T$ = temperature in Kelvin) then $\Delta G$ for given cell reaction at 27$^\circ$C will be\\
$\dfrac12\mathrm{H_2}+\mathrm{AgCl}\rightarrow\mathrm{Ag}+\mathrm{H^+}+\mathrm{Cl^-}$\\
If answer is $xF$ then answer $x$.

\paraq{3}{4}
The metal$-$carbon bond in metal carbonyls posses both $\sigma$ \& $\pi$ character. The metal to ligand bonding creates a synergic effect which strengthens the bond between CO and the metal. Due to synergic effect M$-$C bond length \& C$-$O bond length affected.

\qn{3} Consider $[\mathrm{Co(CO)_4}]^{-}$ as X and $[\mathrm{Mn(CO)_6}]^{+}$ as Y then find how many statements are correct?\\
(i) The C $-$ O bond length in X is more than that in Y\\
(ii) The C $-$ O bond energy in X is more than that in Y\\
(iii) Double bond character in metal to carbon bond in X has more than that in Y\\
(iv) Both X and Y have the same EAN of the central metal atom

\qn{4} Consider the following species\\
$[\mathrm{Ni(CO)_4}]$, $[\mathrm{V(CO)_6}]^{-}$, $\mathrm{CO^+}$, $\mathrm{CO_2}$,\\
$\mathrm{H_2CO_3}$, $\mathrm{CH_3COO^-}$, $[\mathrm{Ni(CO)_3P(CH_3)_3}]$,\\
$[\mathrm{Cr(CO)_6}]$, $[\mathrm{NiPF_3(CO)_3}]$\\
If bond length in CO is `$d$'.\\
The number of species having higher C-O bond length than `$d$' is/are `$X$'\\
The number of species having lower C-O bond length than `$d$' is/are `$Y$'\\
Find the value of $\left(\dfrac{X-Y}{2}\right)$.

\sect{Section-II(ii)}

\qn{5} A KCl solution of conductivity 0.14 S m$^{-1}$ shows a resistance of 4.19 $\Omega$ in a conductivity cell. If the same cell is filled with an HCl solution, the resistance drops to 1.03 $\Omega$. The conductivity of the HCl solution is \underline{\hspace{2em}}$\times10^{-2}$ S m$^{-1}$. (Round off to the Nearest Integer).

\qn{6} A solid cube of edge length $=25.32$ mm of an ionic compound which has NaCl type lattice is added to 1kg of water. The boiling point of this solution is found to be 100.52$^\circ$C (assume 100\% ionisation of ionic compound). If radius of anion of ionic solid is 200 pm then calculate radius of cation of solid in pm (picometer).\\
($k_b$ of water $=0.52$ K kg mole$^{-1}$, Avogadro number, $N_A=6\times10^{23}$, $(\sqrt[3]{75})=4.22$)

\qfullbegin
\qn{7} How many of the following chemical reactions are correct regarding the product given.
\qfig{p2_chem_s4_q07_1}
\qfullsplit
\qfig{p2_chem_s4_q07_2}
\qn{8} How many compounds gives +ve tollen's test and also show mutarotation :
\qfig{p2_chem_s4_q08_1}
\qfullsplit
\qfig{p2_chem_s4_q08_2}
\qfullend

\qn{9} In the given following ores, find out total number of ores which is reduced by carbon reduction method during commercial metallurgical process.\\
Argentite, Chalcopyrite, Galena (lowgrade), Haemetite, Calamine, Casseterite, Cinnabar, Bauxite

\qn{10} Write the sum of the oxidation number of metal in products of given reactions :-\\
$\cdot$ Zn + NaOH$_{\mathrm{aq.}}$ $\rightarrow$\\
$\cdot$ Cu + dil. H$_2$SO$_4$ $\rightarrow$\\
$\cdot$ FeS$_2$ + O$_2$ $_{\mathrm{(Excess)}}$ $\xrightarrow{\Delta}$\\
$\cdot$ Calamine $\xrightarrow{\mathrm{calcination}}$
\qfig{p2_chem_s4_q10_1}

\partsec{3}{MATHEMATICS}
\sect{Section-I(i)}

\qn{1} Coefficient of $x^{12}$ in $(1+x^3)^5(1+x^4)^8(1+x^5)^{11}$ is $-$
\begin{qoptions}
\opt{A}{496}
\opt{B}{501}
\opt{C}{440}
\opt{D}{56}
\end{qoptions}

\qn{2} Let $R$ denotes the set of all non empty relations on the set $A$, where $A=\{1,2,3,4,5\}$. A function $f(x)$ is chosen at random from set $R$. The probability that $f(x)$ is bijective, $f(x)\ne x\ \forall\ x\in A$ and $f(1)\ne2$, is $-$
\begin{qoptions}
\opt{A}{$\dfrac{33}{2^{25}}$}
\opt{B}{$\dfrac{33}{2^{25}-1}$}
\opt{C}{$\dfrac{27}{2^{25}}$}
\opt{D}{$\dfrac{27}{2^{25}-1}$}
\end{qoptions}

\qn{3} Greatest term in the expansion of $(7-4x)^{10}$ is (where $x=1$)
\begin{qoptions}
\opt{A}{$T_4$ \& $T_5$}
\opt{B}{$T_4$ only}
\opt{C}{$T_5$ only}
\opt{D}{$T_5$ \& $T_6$}
\end{qoptions}

\qn{4} For three non impossible events $A$, $B$ and $C$, $P(A\cap B\cap C)=0$, $P(A\cup B\cup C)=\dfrac34$, $P(A\cap B)=\dfrac13$ and $P(C)=\dfrac16$. The probability, exactly one of $A$ or $B$ occurs but $C$ doesn't occur is $-$
\begin{qoptions}
\opt{A}{$\dfrac1{12}$}
\opt{B}{$\dfrac56$}
\opt{C}{$\dfrac14$}
\opt{D}{$\dfrac23$}
\end{qoptions}

\sect{Section-I(ii)}

\qn{5} If $A,B,C$ are three independent \& exhaustive events such that $P\left((A\cap\overline B)\cup(\overline A\cap B)\right)=\dfrac12$, $P\left(A\cap B\cap\overline C\right)=\dfrac7{15}$, then $P(A\cap B)$ less than
\begin{qoptions}
\opt{A}{$\dfrac12$}
\opt{B}{$\dfrac23$}
\opt{C}{$\dfrac35$}
\opt{D}{$\dfrac13$}
\end{qoptions}

\qn{6} A fair die is rolled `$n$' times. If $P(n)$ denotes the probability that there are atleast two equal numbers among the result obtained then $-$
\begin{qoptions}
\opt{A}{$P(n\geqslant7)=1$}
\opt{B}{$P(n=2)=\dfrac16$}
\opt{C}{$P(n=4)=\dfrac{13}{18}$}
\opt{D}{$P(n=6)=\dfrac{317}{324}$}
\end{qoptions}

\qn{7} For the equation $2\cos2x\left(\cos2x-\cos\left(\dfrac{2014\pi^2}{x}\right)\right)$ $=\cos4x-1$ ; which of the following is/are correct ?
\begin{qoptions}
\opt{A}{The number of solutions of the equation is 3.}
\opt{B}{The sum of all the positive solutions is $1080\pi$}
\opt{C}{The number of positive solutions is 4}
\opt{D}{The sum of all positive solutions is $1008\pi$}
\end{qoptions}

\sect{Section-II(i)}

\paraq{1}{2}
Let bag `A' contains 2 Red \& 3 Green balls \& bag `B' contains 3 Red \& 2 Green balls. If two balls are drawn randomly from each of bag A \& B and then put in empty bag C.

\qn{1} Probability that bag C contains 3 Red balls, is $-$

\qn{2} If it is found that bag C contains equal number of Red and Green balls, then probability that 2 Red balls were drawn from bag B, is $-$

\paraq{3}{4}
Consider the word $W=$ TALLENTEX.

\qn{3} The number of permutations that can be formed out of the letters of the word TALLENTEX taking three letters at a time is $-$

\qn{4} The number of rearrangements of the letters of the word $W$ such that ALLEN is always formed is $-$

\sect{Section-II(ii)}

\qn{5} If $1\cdot\dbinom{2014}{1}\dbinom{2013}{2013}+2\dbinom{2014}{2}\dbinom{2013}{2012}+3\dbinom{2014}{3}\dbinom{2013}{2011}+\ldots 2014\dbinom{2014}{2014}\dbinom{2013}{0}=2014\left({}^{4026}C_\lambda\right)$, then sum of all the digits of $\lambda$ is

\qn{6} The number of solution satisfying the equation $\begin{vmatrix}1+\sin^3\theta&1\\-\cos^3\theta&1\end{vmatrix}=\begin{vmatrix}\cos\theta&4\cos\theta\\0&3\sin\theta\end{vmatrix}$ in $\theta\in(0,4\pi)$ is equal to

\qn{7} Four distinct integer are picked at random from $\{0,1,2,3,4,5,6\}$. If the probability that among those selected the second smallest is 3 is `$P$', then `$35P$' is equal to

\qn{8} Let $x_1,x_2,\ldots,x_{100}$ are 100 observations such that $\sum x_i=0$, $\displaystyle\sum_{1\leqslant i<j\leqslant100}\left|x_ix_j\right|=80000$ \& mean deviation from their mean is 5, then their standard deviation, is $-$

\qn{9} If $A=\dfrac{\log_{18}21\cdot\log_{23}52\cdot\log_{21}51\cdot\log_{52}18}{\log_{11}51\cdot\log_{23}11}$ and\\
$B=\log_9(\tan^226^\circ)+\log_9(\tan^228^\circ)+\log_9(\tan^260^\circ)+\log_9(\tan^262^\circ)+\log_9(\tan^264^\circ)$, then $\dfrac AB$ is

\qn{10} Consider a grid as shown in figure having 8 blocks of dimension $1\times1$. Each block is to be painted with either blue or red colour. In how many ways this can be done if atleast one block of dimension $2\times2$ is to be painted red completely ?
\qfig{p2_math_s4_q10_1}

\end{multicols}
\end{document}
